\documentclass[]{article}

\usepackage{amsmath}
\usepackage{multirow}
\usepackage{natbib}
\usepackage{graphicx} 


\begin{document}

\subsection{Cosmological likelihood from DESI DR2}
The foundation of our Bayesian evidence calculation is the robust extraction of the cosmological likelihood $P(D_{\text{cosmo}} | \Sigma m_\nu)$ from the DESI DR2 MCMC chains. For the baseline $\Lambda$CDM model, utilizing the DESI BAO and Planck CamSpec CMB data, the marginalized 1D posterior yields a 95\% upper limit of $\Sigma m_\nu < 0.0642$\,eV. To continuously evaluate the evidence integrals, we modeled this posterior using a truncated Gaussian distribution constrained to the physical region $\Sigma m_\nu \geq 0$. As shown in Figure~\ref{fig:image0}, the optimal fit for the baseline dataset resulted in a central value of $\mu_0 = -0.009$\,eV and a standard deviation of $\sigma = 0.036$\,eV. This representation accurately captures the bulk of the posterior probability while naturally accommodating the physical boundary at zero mass.

The validity of our analytical approximation is confirmed by its close agreement with the Feldman-Cousins profile likelihood reported by the DESI collaboration ($\mu_0 = -0.036$\,eV, $\sigma = 0.043$\,eV), also shown in Figure~\ref{fig:image0}. Furthermore, as illustrated in Figure~\ref{fig:image1}, alternative CMB likelihoods, such as plik (95\% UL = 0.0691\,eV) and L-H (95\% UL = 0.0774\,eV), demonstrate consistent, albeit slightly relaxed, constraints. The stability of the likelihood across these different CMB treatments ensures that our subsequent evidence calculations are not overly sensitive to specific CMB pipeline choices.

\begin{figure}[h!]
    \centering
    \includegraphics[width=\columnwidth]{../input_files/plots/likelihood_comparison_1_20260501_180057.png}
    \caption{The marginalized one-dimensional posterior probability density for the sum of neutrino masses, $\Sigma m_\nu$, derived from DESI DR2 and Planck CamSpec data under a $\Lambda$CDM model. The posterior, estimated via Kernel Density Estimation (KDE) from the MCMC chains (solid black), is modeled using a truncated Gaussian distribution (dashed blue; $\mu_0 = -0.009$ eV, $\sigma = 0.036$ eV) for Bayesian evidence calculations. The close agreement with the Feldman-Cousins profile likelihood (dash-dotted red; $\mu_0 = -0.036$ eV, $\sigma = 0.043$ eV) validates this analytical approximation. The analysis yields a 95\% upper limit of $\Sigma m_\nu < 0.0642$ eV (dotted vertical line).}
    \label{fig:image0}
\end{figure}

\begin{figure}[h!]
    \centering
    \includegraphics[width=\columnwidth]{../input_files/plots/likelihood_all_datasets_1_20260501_180057.png}
    \caption{Marginalized posterior distributions for the sum of neutrino masses, $\Sigma m_\nu$, derived from DESI DR2 data combined with different Planck CMB likelihoods (CamSpec, plik, and L-H) within the $\Lambda$CDM model. The constraints are consistent across the different CMB treatments, demonstrating the robustness of the likelihood used for the Bayesian evidence calculation, with the baseline CamSpec likelihood providing the tightest constraint (95\% upper limit of 0.0642 eV). All posteriors show that the probability density is concentrated at values significantly below the minimum mass required for the Inverted Hierarchy (IH, 0.099 eV) while being consistent with the Normal Hierarchy (NH, 0.059 eV).}
    \label{fig:image1}
\end{figure}

\subsection{Bayesian evidence and Bayes factors}
The core quantitative result of this study is the Bayes factor $K = P(D|\text{NH}) / P(D|\text{IH})$, which measures the relative probability of the Normal Hierarchy (NH) over the Inverted Hierarchy (IH) given the combined cosmological and oscillation data. We computed $K$ under two fundamentally distinct prior frameworks: the physically motivated hierarchical SJPV prior and the objective, information-theoretic HS prior.

For the baseline DESI DR2 dataset combined with NuFIT 6.0 oscillation parameters, the evidence is overwhelmingly in favor of the Normal Hierarchy. Under the SJPV prior, the base Bayes factor—derived solely from the cosmological likelihood and prior volume, explicitly excluding the $\Delta\chi^2$ preference from the oscillation global fit—is $K_{\text{base}} = 484.5$. When the $\Delta\chi^2 = 6.1$ preference from the NuFIT 6.0 global fit is multiplicatively included ($K_{\text{full}} = K_{\text{base}} \times \exp(6.1/2)$), the total Bayes factor reaches a staggering $K_{\text{full}} = 10231.4$. On the Jeffreys scale, this constitutes ``decisive'' evidence ($K > 100$) against the Inverted Hierarchy.

Crucially, this decisive preference is not merely an artifact of the hierarchical prior's geometric properties. When employing the objective HS prior—which imposes no geometric penalty on the degenerate mass spectrum required by the IH—the base Bayes factor is $K_{\text{base}} = 22.0$. Including the oscillation preference yields a total Bayes factor of $K_{\text{full}} = 464.6$. The fact that $K_{\text{full}}$ exceeds the decisive threshold of 100 even under this conservative, agnostic prior demonstrates that the DESI DR2 data alone is sufficiently powerful to rule out the Inverted Hierarchy. These results, along with those from alternative datasets, consistently yield decisive Bayes factors and are summarized comprehensively in Figure~\ref{fig:image8}.

\begin{figure}[h!]
    \centering
    \includegraphics[width=\columnwidth]{../input_files/plots/step_7_bayes_factors_summary_7_20260501_182746.png}
    \caption{Comparison of the total Bayes factor, K, supporting the Normal Hierarchy over the Inverted Hierarchy across various cosmological datasets. Results are shown for the physically motivated SJPV prior (blue) and the objective HS prior (red). For the baseline $\Lambda$CDM model, incorporating different CMB likelihoods or supernovae datasets, both priors yield decisive evidence (K > 100) for the Normal Hierarchy. The preference remains robust even within an extended $w_0w_a$CDM dark energy model, where the evidence is still decisive under the SJPV prior and strong under the HS prior.}
    \label{fig:image8}
\end{figure}

\subsection{Historical evolution of the evidence}
To contextualize the strength of the DESI DR2 constraints, we reconstructed the historical evolution of the Bayes factor from 2002 to the present. As illustrated in Figure~\ref{fig:image2}, the evidence for the NH has grown monotonically as cosmological upper limits on $\Sigma m_\nu$ have progressively tightened. In the early 2000s (e.g., 2dFGRS, $\Sigma m_\nu < 1.80$\,eV), the Bayes factor was near unity, indicating no statistical preference. By the Planck 2015 era ($\Sigma m_\nu < 0.18$\,eV), the SJPV prior began to show ``strong'' evidence ($K > 10$), driven primarily by the geometric volume effect as the allowed parameter space shrank toward the IH minimum mass threshold ($\Sigma_{\text{min}}^{\text{IH}} \approx 0.099$\,eV). However, under the objective HS prior, the evidence remained weak.

The transition to the DESI DR2 era ($\Sigma m_\nu < 0.0642$\,eV) marks a critical inflection point: the cosmological upper limit has now plunged significantly below the IH minimum mass. Consequently, the Bayes factor under both prior frameworks has surged past the ``decisive'' threshold ($K > 100$). This historical trajectory validates previous methodologies that foresaw the impending collapse of the IH parameter space and visually underscores the transformative, conclusive impact of the DESI DR2 data.

\begin{figure}[h!]
    \centering
    \includegraphics[width=\columnwidth]{../input_files/plots/step_3_historical_bayes_factor_3_20260501_181318.png}
    \caption{Historical evolution of the Bayes factor K, quantifying the evidence for the Normal Hierarchy (NH) versus the Inverted Hierarchy (IH), from 2002 to 2025. The analysis is presented for two distinct prior frameworks: the physically motivated SJPV prior (blue circles) and the objective HS prior (red squares). The plot illustrates a monotonic growth in evidence for the NH as cosmological constraints on the sum of neutrino masses ($\Sigma m_\nu$) have progressively tightened over time. The most recent data point, corresponding to the DESI DR2 era, represents a critical inflection point where the cosmological upper limit on $\Sigma m_\nu$ has fallen below the minimum mass required by the IH. This has caused the Bayes factor under both prior frameworks to surge past the ``decisive'' threshold (K > 100), demonstrating that the preference for the NH is a robust, data-driven conclusion.}
    \label{fig:image2}
\end{figure}

\subsection{Prior sensitivity and geometric volume effects}
The difference in magnitude between the SJPV and HS Bayes factors ($K_{\text{full}} \approx 10231$ vs $464$) highlights the role of the prior in Bayesian model comparison. The SJPV prior, assuming exchangeable masses from a common log-normal distribution, induces a profound geometric volume penalty on the IH. The IH requires two nearly degenerate heavy masses ($m_1 \simeq m_2 \gg m_3$), a configuration that is statistically less probable in the hyperparameter space than the NH configuration. As visualized in Figure~\ref{fig:image3}, the accessible parameter space for the IH under the SJPV prior is severely restricted compared to the NH. Sensitivity checks, where we expanded the hyperprior bounds by orders of magnitude, resulted in negligible changes to the Bayes factor, confirming the stability of the hierarchical framework.

Conversely, the HS prior is constructed to be minimally informative. Derived from the Fisher information matrix of the oscillation observables, the resulting prior, $P_{\text{HS}} \propto m_L m_M + m_L m_H + m_M m_H$, acts as an objective reference that removes the geometric penalty associated with exchangeability. The convergence of both priors to a ``decisive'' conclusion with DESI DR2 data confirms that while the SJPV prior accelerates the statistical preference through physical arguments, the underlying cosmological data is now strong enough to independently drive the result.

\begin{figure}[h!]
    \centering
    \includegraphics[width=\columnwidth]{../input_files/plots/step_4_prior_volume_sensitivity_4_20260501_181745.png}
    \caption{Comparison of prior probability densities for the sum of neutrino masses ($\Sigma m_\nu$) under the SJPV (left) and HS (middle) frameworks. The SJPV prior exhibits a strong geometric volume penalty, restricting the parameter space for the Inverted Hierarchy (IH, red dashed) relative to the Normal Hierarchy (NH, blue solid). The objective HS prior removes this penalty, assigning comparable volumes. The right panel confirms the stability of the decisive Bayes factor for the NH under the SJPV prior, showing it is insensitive to the expansion of its hyperparameter bounds.}
    \label{fig:image3}
\end{figure}

\subsection{Individual mass posteriors and spectra}
The decisive preference for the NH allows us to construct highly constrained posterior distributions for the individual neutrino mass eigenstates. For the NH, the posterior shown in Figure~\ref{fig:image4} is strongly peaked near the minimum allowed mass configuration. The lightest state $m_1$ is constrained to be nearly massless ($< 0.01$\,eV), and consequently, the masses of $m_2$ and $m_3$ are dictated almost entirely by the oscillation splittings, yielding highly localized posteriors at $m_2 \approx 0.008$\,eV and $m_3 \approx 0.050$\,eV. The sum $\Sigma m_\nu$ tightly hugs its theoretical minimum of $\sim 0.059$\,eV, in perfect agreement with the DESI DR2 constraints.

In stark contrast, the posterior for the IH (Figure~\ref{fig:image5}) is confined to $\Sigma m_\nu \gtrsim 0.099$\,eV. This entire parameter space is in severe tension with the cosmological likelihood. The mass spectrum plot in Figure~\ref{fig:image6} visually reinforces this conclusion. The NH curve naturally accommodates a $\Sigma m_\nu$ well within the DESI DR2 95\% upper limit, whereas the IH curve shows that even as the lightest mass approaches zero, the sum $\Sigma m_\nu$ cannot drop below $\sim 0.099$\,eV, placing the entire hierarchy in conflict with observations.

\begin{figure}[h!]
    \centering
    \includegraphics[width=\columnwidth]{../input_files/plots/step_5_corner_plot_NH_5_20260501_182112.png}
    \caption{Joint and marginal posterior distributions for the individual neutrino masses ($m_1$, $m_2$, $m_3$) and their sum ($\Sigma m_\nu$) for the Normal Hierarchy (NH), derived from the baseline DESI DR2 data combined with the SJPV prior. The posterior is strongly peaked near the minimum allowed mass configuration, showing the lightest mass state, $m_1$, to be nearly massless. Consequently, the masses of $m_2$ and $m_3$ are sharply determined by oscillation splittings to be approximately 0.008 eV and 0.050 eV, respectively. This forces the total mass, $\Sigma m_\nu$, to be tightly constrained around its theoretical minimum of $\sim 0.059$ eV, in full agreement with the cosmological constraints.}
    \label{fig:image4}
\end{figure}

\begin{figure}[h!]
    \centering
    \includegraphics[width=\columnwidth]{../input_files/plots/step_5_corner_plot_IH_5_20260501_182113.png}
    \caption{Joint and marginal posterior distributions for the individual neutrino masses ($m_1, m_2, m_3$) and their sum ($\Sigma m_\nu$) under the Inverted Hierarchy (IH) model, using the baseline DESI DR2 data and the SJPV prior. The plot illustrates the required mass degeneracy ($m_1 \approx m_2 \approx 0.05$ eV) and shows that the posterior for the total mass is confined above the IH minimum of $\sim 0.099$ eV. This entire parameter space is in severe tension with the 95\% cosmological upper limit of $0.0642$ eV, visually demonstrating why the IH is decisively disfavored by the data.}
    \label{fig:image5}
\end{figure}

\begin{figure}[h!]
    \centering
    \includegraphics[width=\columnwidth]{../input_files/plots/step_5_mass_spectrum_5_20260501_182114.png}
    \caption{Neutrino mass spectrum for the Normal (NH) and Inverted (IH) Hierarchies, plotting individual masses and their sum ($\Sigma m_\nu$) as a function of the lightest neutrino mass ($m_L$). The plot visually demonstrates the tension between the IH and cosmological data. The minimum possible sum for the NH ($\sim 0.059$ eV) is well within the DESI DR2 95\% upper limit of 0.0642 eV, whereas the minimum for the IH ($\sim 0.099$ eV) lies significantly above it, placing the entire IH parameter space in conflict with observations.}
    \label{fig:image6}
\end{figure}

\subsection{Implications for neutrinoless double-beta decay}
The resolution of the mass hierarchy has profound implications for the search for neutrinoless double-beta decay ($0\nu\beta\beta$). The decay rate is governed by the effective Majorana mass, $m_{\beta\beta}$. By sampling over the PMNS mixing parameters and unknown Majorana CP phases, we computed the posterior distribution for $m_{\beta\beta}$ under both hierarchies, shown in Figure~\ref{fig:image7}.

For the favored NH, the posterior for $m_{\beta\beta}$ is heavily suppressed due to the near-zero value of $m_1$ and potential phase cancellations. We find a 95\% credible interval of $[0.95, 11.55]$\,meV, with a median of $3.28$\,meV. This range lies almost entirely below the projected sensitivity of next-generation experiments like nEXO ($\sim 5$\,meV). In contrast, the now-disfavored IH would have predicted a median $m_{\beta\beta}$ of $37.03$\,meV (95\% CI: $[18.36, 49.51]$\,meV), well within experimental reach. Our results suggest that $0\nu\beta\beta$ experiments face a significantly more challenging landscape, with the implied lower bound on the half-life $T^{0\nu\beta\beta}_{1/2}$ for the NH extending well beyond $10^{28}$ years.

\begin{figure}[h!]
    \centering
    \includegraphics[width=\columnwidth]{../input_files/plots/step_6_mbb_projections_6_20260501_182338.png}
    \caption{Posterior probability distributions for the effective Majorana mass ($m_{\beta\beta}$) for the Normal Hierarchy (NH, blue) and Inverted Hierarchy (IH, red). The left panel shows the NH posterior is concentrated at low values, with a 95\% credible interval of [0.95, 11.55] meV, largely below the projected sensitivity of next-generation experiments like nEXO. Conversely, the IH posterior predicts a much larger signal with a 95\% credible interval of [18.36, 49.51] meV, which would be well within experimental reach. The right panel translates these distributions into the implied lower bound on the neutrinoless double-beta decay half-life ($T_{1/2}$), demonstrating that the NH scenario points towards a half-life likely exceeding $10^{28}$ years, making a non-detection in upcoming experiments the most probable outcome.}
    \label{fig:image7}
\end{figure}

\subsection{Robustness to the dark energy model}
A critical test of our conclusion is its robustness against extensions to the standard cosmological model. We re-evaluated the Bayesian evidence using a $w_0w_a$CDM model, where a time-varying dark energy equation of state can be degenerate with $\Sigma m_\nu$. In this extended parameter space, the 95\% upper limit on $\Sigma m_\nu$ relaxes to $0.163$\,eV.

As shown in Figure~\ref{fig:image8}, the Bayes factors decrease as the tension with the IH minimum mass is alleviated. The SJPV prior yields $K_{\text{full}} = 482.1$, maintaining a ``decisive'' classification. The objective HS prior drops to $K_{\text{full}} = 42.6$. While this is a significant reduction from the $\Lambda$CDM baseline ($K_{\text{full}} = 464.6$), a Bayes factor of 42.6 still constitutes ``strong'' evidence for the Normal Hierarchy on the Jeffreys scale. This demonstrates that the statistical preference for the NH remains robust and highly significant even when allowing for a dynamic dark energy equation of state.

\end{document}
                